\section{Conclusion}
\label{sec:conclusion}

We constructed a 2x2 ablation of data quality and model capacity in from-scratch Transformer machine translation, and found that the capacity return changes sign with data quality: $+0.56$ test BLEU on the strict-filter 9.3M en-fr corpus, $-0.87$ on the loose-filter 30M corpus drawn from the same source. The same architectural change pays in one regime and costs in another. The data-quality dimension qualitatively replicates on en-de at 4.5M scale (Base $>$ Big by 1.57 BLEU), supporting an interpretation under which Big requires both clean data \emph{and} enough of it; failing either condition flips the sign of the capacity return.

A natural follow-up --- can we recover v2's noise via post-training quality filtering? --- yields a clean negative result. Reference-free QE-filtered fine-tuning (CometKiwi-22, top-1M of the 30M v2 corpus, 6K fine-tuning steps under cross-entropy loss) degrades newstest BLEU by 1.5--2.3 on both Base and Big v1.1, with Big losing more. The mechanism is distinct from the noise-amplification of \S\ref{sec:enfr_v2}: the fine-tuning data is high-quality but heavily UN/legislative, mismatched to news domain. We argue that this decomposes ``data quality'' into translation correctness and target-distribution alignment --- two independent axes, of which reference-free QE measures only the first.

The paper's prescription, in three steps:
\begin{enumerate}
\item \textbf{Filter for noise.} Cap data at the strict-cleaning frontier; do not chase scale past the noise wall.
\item \textbf{Match the eval distribution.} QE-filtered fine-tuning or distillation \emph{must} condition on domain or register; absolute QE on its own is insufficient and can hurt.
\item \textbf{Then scale capacity.} The capacity return is $+0.56$ when (1) and (2) hold, and is negative or near zero when either fails.
\end{enumerate}

We close with a methodology caveat about loss-spike-guard trigger rates as a data-noise diagnostic: trigger rates are joint functions of data noise and model convergence depth, and only same-capacity / different-data comparisons are clean. Cross-capacity comparisons that we relied on in early drafts of this work were retroactively invalidated by this realization; the corrected analysis is in \S\ref{sec:methodology}.

\paragraph{What we did not test.}
Only two language pairs (en-fr, en-de). Single seed per configuration. No human evaluation. No back-translation augmentation. No domain-conditional QE filter --- the natural fix for the fine-tuning failure of \S\ref{sec:sft}, and the most concrete future-work direction. All checkpoints, code, configs, and stdout logs are released to enable independent replication.

\paragraph{What we hope is taken from this paper.}
The headline empirical claim: one cannot read off the value of extra capacity without knowing what data it will see, and one cannot read off the value of more data without knowing whether the post-training pipeline will move the model in the right direction. ``More data'' and ``more parameters'' are conditioned on each other and on the post-training pipeline. The 2x2 design we used --- {strict, loose} x {Base, Big}, on a shared source pool, with cross-language replication and a closed-loop fine-tuning test --- is itself a transferable template. We hope future scaling-and-data papers report it more often than they currently do.
